\chapter{Attention and the Transformer}
\label{ch:transformer}

The Transformer \cite{vaswani} is the architecture behind essentially every large
language model and much of modern vision and speech. It is also the most common
``implement it from scratch'' request in ML interviews, from founder rounds at
start-ups to research-engineer loops. This chapter builds it piece by piece with the
shapes written out at every step, derives why the scores are divided by $\sqrt{d_k}$,
counts parameters and FLOPs, and ends with a capstone: a complete encoder--decoder
Transformer written in PyTorch and again in NumPy, trained on a toy task, with the two
implementations and PyTorch's own layers asserted to agree.

\section{Attention as a soft dictionary lookup}
\prototype{A decoder translating word $t$ looks back at all encoder states and takes a
weighted average, with weights from a softmax over relevance scores.}

A Python dictionary returns the value whose key \emph{equals} the query. Attention returns
a blend of \emph{all} values, weighted by how well each key matches the query:
\[
	\text{output}=\sum_i\alpha_i\,\vv_i,\qquad
	\alpha_i=\frac{\exp(\mathrm{score}(\bq,\bk_i))}{\sum_j\exp(\mathrm{score}(\bq,\bk_j))}.
\]
Because the weights come from a softmax, the lookup is differentiable and can be learned.

Attention entered deep learning to fix the seq2seq bottleneck of \cref{ch:rnn}: instead of
squeezing the source sentence into one vector, the decoder at step $t$ computes a fresh
\vocab{context vector} $\bc_t=\sum_i\alpha_{ti}\h_i$ over all encoder states $\h_i$.
\begin{itemize}
	\ii \vocab{Additive (Bahdanau) attention} \cite{bahdanau} scores with a one-hidden-layer
	network on the previous decoder state and an encoder state,
	$\mathrm{score}(\bs_{t-1},\h_i)=\vv\T\tanh(\W\bs_{t-1}+\U\h_i)$.
	\ii \vocab{Multiplicative (Luong) attention}: $\mathrm{score}(\bs_t,\h_i)=\bs_t\T\W\h_i$ (or just
	the dot product $\bs_t\T\h_i$). Cheaper: one matrix multiplication for all pairs.
\end{itemize}
The Transformer keeps the dot product, adds a scale, and, crucially, uses attention between
positions of the \emph{same} sequence (\vocab{self-attention}) as its main operation, dropping
recurrence entirely.

\section{Scaled dot-product attention}
\prototype{Three tokens with $2$-dimensional queries, keys and values; the first output row is
$(\gv{c08_attention/o1_0},\,4.000)$.}

\subsection{Definition and a worked example}
Stack queries, keys and values as rows: $\Q\in\RR^{T_q\times d_k}$, $\K\in\RR^{T_k\times d_k}$,
$\V\in\RR^{T_k\times d_v}$. Then
\[
	\Attn(\Q,\K,\V)=\softmax\!\Bigl(\frac{\Q\K\T}{\sqrt{d_k}}\Bigr)\V,
\]
with the softmax taken over each row (over keys). The score matrix has shape $\shp{T_q,T_k}$;
each row of the weight matrix sums to one; the output has shape $\shp{T_q,d_v}$.

\codefile[\script{c08_attention}]{gen/dl/snip/c08_attention-sdpa.py}

\begin{example}[Attention on three tokens]
	\label{ex:attn3}
	$\Q=\K=\begin{psmallmatrix}1&0\\0&1\\1&1\end{psmallmatrix}$,
	$\V=\begin{psmallmatrix}1&2\\3&4\\5&6\end{psmallmatrix}$, $d_k=2$.
	\[
		\Q\K\T=\begin{pmatrix}\gv{c08_attention/QKt}\end{pmatrix},\quad
		\frac{\Q\K\T}{\sqrt2}=\begin{pmatrix}\gv{c08_attention/S}\end{pmatrix},
	\]
	\[
		\A=\begin{pmatrix}\gv{c08_attention/A}\end{pmatrix},\quad
		\A\V=\begin{pmatrix}\gv{c08_attention/O}\end{pmatrix}.
	\]
	Token 3 matches every key and attends most to itself; token 1 splits its attention equally
	between keys 1 and 3. The result equals \pt{F.scaled\_dot\_product\_attention}.
	\begin{center}
	\begin{tikzpicture}
	\begin{axis}[width=0.36\linewidth, height=0.36\linewidth, enlargelimits=false, axis on top,
		xlabel={key}, ylabel={query}, y dir=reverse, xtick={0,1,2}, ytick={0,1,2},
		xticklabels={1,2,3}, yticklabels={1,2,3}, colormap/viridis, colorbar,
		colorbar style={width=0.25cm, font=\scriptsize}, point meta min=0, point meta max=1, font=\small]
		\addplot[matrix plot*, point meta=explicit, mesh/cols=3, mesh/ordering=y varies]
			table[x=x, y=y, meta=a] {gen/dl/data/c08_attention-heat3.dat};
	\end{axis}
	\end{tikzpicture}
	\end{center}
\end{example}

\subsection{Why divide by $\sqrt{d_k}$}
\begin{proposition}
	If the components of $\bq$ and $\bk$ are independent with mean $0$ and variance $1$, then
	$\bq\T\bk=\sum_{i=1}^{d_k}q_ik_i$ has mean $0$ and variance $d_k$.
\end{proposition}
\begin{proof}
	Each term has mean $\E q_i\E k_i=0$ and variance $\E q_i^2\E k_i^2=1$; the $d_k$ terms are
	independent, so the variances add.
\end{proof}
So unscaled scores have standard deviation $\sqrt{d_k}$, and for $d_k=64$ typical score gaps are
around $8$: the softmax is then nearly one-hot. A saturated softmax has a Jacobian
$\diag(\bm p)-\bm p\bm p\T$ close to zero, so the gradients through the attention weights vanish and
learning stalls. Dividing by $\sqrt{d_k}$ restores unit variance whatever the head size.

\begin{example}[Measuring it]
	Random unit-variance $\bq,\bk$; variance of $\bq\T\bk$: $\gv{c08_attention/var_4}$,
	$\gv{c08_attention/var_16}$, $\gv{c08_attention/var_64}$, $\gv{c08_attention/var_256}$,
	$\gv{c08_attention/var_1024}$ for $d_k=4,16,64,256,1024$. Mean largest attention weight over
	$32$ keys, without and with scaling:
	\begin{center}
		\small
		\begin{tabular}{lccccc}
			\toprule
			$d_k$ & 4 & 16 & 64 & 256 & 1024 \\
			\midrule
			unscaled & $\gv{c08_attention/maxraw_4}$ & $\gv{c08_attention/maxraw_16}$ & $\gv{c08_attention/maxraw_64}$ & $\gv{c08_attention/maxraw_256}$ & $\gv{c08_attention/maxraw_1024}$ \\
			scaled & $\gv{c08_attention/maxsc_4}$ & $\gv{c08_attention/maxsc_16}$ & $\gv{c08_attention/maxsc_64}$ & $\gv{c08_attention/maxsc_256}$ & $\gv{c08_attention/maxsc_1024}$ \\
			\bottomrule
		\end{tabular}
	\end{center}
	At $d_k=1024$ the unscaled weights have entropy $\gv{c08_attention/ent_raw_1024}$ nats against
	$\gv{c08_attention/ent_sc_1024}$ scaled (uniform: $\ln32=\gv{c08_attention/ln32}$). For one row
	of $32$ scores at $d_k=512$ the Frobenius norm of the softmax Jacobian is
	$\gv{c08_attention/jac_raw}$ unscaled and $\gv{c08_attention/jac_scaled}$ scaled.
	\begin{center}
	\begin{tikzpicture}
	\begin{axis}[napkinplot, width=0.55\linewidth, height=0.36\linewidth, xmode=log, log basis x=2,
		xlabel={$d_k$}, ylabel={mean max weight}, ymin=0, ymax=1, legend pos=outer north east]
		\addplot+[mark=*] table[x=dk,y=max_raw] {gen/dl/data/c08_attention-sqrtdk.dat}; \addlegendentry{$\Q\K\T$}
		\addplot+[mark=square*] table[x=dk,y=max_scaled] {gen/dl/data/c08_attention-sqrtdk.dat}; \addlegendentry{$\Q\K\T/\sqrt{d_k}$}
	\end{axis}
	\end{tikzpicture}
	\end{center}
\end{example}

\begin{remark}
	\trap ``We divide by $\sqrt{d_k}$ to normalise the attention weights'' (the softmax already does
	that) or ``to reduce computation''. The reason is the \emph{variance} of the dot product: it keeps
	the softmax out of saturation so gradients flow. It is $\sqrt{d_k}$, not $d_k$, because standard
	deviation scales with the square root of the variance.
\end{remark}

\subsection{Masks}
A mask sets blocked scores to $-\infty$ \emph{before} the softmax, so blocked positions get weight
exactly $0$ and each row still sums to one.
\begin{itemize}
	\ii \vocab{Causal (look-ahead) mask}: query $i$ may attend only to keys $j\le i$ (strictly upper
	triangle blocked). Needed in decoders so that position $i$ cannot see the token it is trained to
	predict; it is what lets a whole sequence be trained in parallel with teacher forcing.
	\ii \vocab{Padding mask}: in a batch of sequences of different lengths, padding positions are
	blocked as keys so no real token attends to them.
\end{itemize}
For \cref{ex:attn3}, a causal mask and, alternatively, treating key 3 as padding give the weights
\[
	\A_\text{causal}=\begin{pmatrix}\gv{c08_attention/Ac}\end{pmatrix},\qquad
	\A_\text{pad}=\begin{pmatrix}\gv{c08_attention/Ap}\end{pmatrix},
\]
matching PyTorch with \pt{is\_causal=True} and with an \pt{attn\_mask} respectively.

\begin{remark}
	\trap Masking \emph{after} the softmax (multiplying the weights by $0/1$) leaves rows that no longer
	sum to one: here the row sums would be $\gv{c08_attention/Abad_rowsum}$. And the causal mask is not
	optional at inference ``because future tokens are not there yet'': during training they are there,
	and the model must be trained under the same constraint it will be used with.
\end{remark}

\section{Multi-head attention, with shapes}
\prototype{$\dmodel=512$, $8$ heads of size $64$: four $512\times512$ projections, $\gv{c08_attention/mha_512}$
parameters, and attention computed in eight $64$-dimensional subspaces in parallel.}

One attention pattern per layer is limiting: a token may need to look at its syntactic head, its
coreferent and the previous token at the same time. \vocab{Multi-head attention} runs $h$ attentions in
parallel on learned projections of size $d_k=\dmodel/h$ and concatenates the results:
\begin{gather*}
	\MHA(\X)=\mathrm{Concat}(\text{head}_1,\dots,\text{head}_h)\,\Wo,\\
	\text{head}_j=\Attn\bigl(\X\Wq^{(j)},\ \X\Wk^{(j)},\ \X\Wv^{(j)}\bigr).
\end{gather*}
In code all heads come from one $\dmodel\times\dmodel$ matrix per role, reshaped:
\begin{center}
	\small
	\begin{tabular}{ll}
		\toprule
		step & shape \\
		\midrule
		input $\X$ & $\shp{B,T,\dmodel}$ \\
		$\X\Wq\T$, $\X\Wk\T$, $\X\Wv\T$ & $\shp{B,T,\dmodel}$ each \\
		split heads: reshape, transpose & $\shp{B,h,T,d_k}$ \\
		scores $\Q\K\T/\sqrt{d_k}$ (+ mask) & $\shp{B,h,T,T}$ \\
		softmax over the last axis, times $\V$ & $\shp{B,h,T,d_k}$ \\
		merge heads: transpose, reshape & $\shp{B,T,\dmodel}$ \\
		output projection $\Wo$ & $\shp{B,T,\dmodel}$ \\
		\bottomrule
	\end{tabular}
\end{center}

\codefile[\script{c08_attention}]{gen/dl/snip/c08_attention-mha.py}

The script copies the weights of a PyTorch \pt{nn.MultiheadAttention} into this function and asserts
that outputs and per-head attention weights agree, with and without a causal mask. Parameter count:
$4\dmodel^2+4\dmodel$ (four projections with biases), $\gv{c08_attention/mha_params}$ for $\dmodel=8$;
$\gv{c08_attention/mha_512}$, $\gv{c08_attention/mha_768}$, $\gv{c08_attention/mha_1024}$ for
$\dmodel=512,768,1024$ (all asserted). The count does not depend on $h$: more heads means smaller heads,
not more parameters.

\begin{remark}
	\trap Self-attention without positional information is \emph{permutation-equivariant}: shuffle the
	input tokens and the outputs are shuffled the same way (asserted in the script). ``The dog bit the
	man'' and ``the man bit the dog'' would produce the same set of vectors. That is why positional
	encodings are not optional.
\end{remark}

\section{Positional encodings}
\prototype{Sinusoidal encodings for $\dmodel=8$: position 1 is
$(\gv{c08_attention/pe_row1})$.}

\begin{itemize}
	\ii \vocab{Sinusoidal} (original Transformer): added to the token embeddings,
	\[ PE_{p,2i}=\sin\bigl(p/10000^{2i/\dmodel}\bigr),\qquad PE_{p,2i+1}=\cos\bigl(p/10000^{2i/\dmodel}\bigr). \]
	Each pair of dimensions is a clock at its own frequency, from fast (one cycle per $2\pi$ positions) to
	very slow. For every offset $k$ there is a fixed rotation $R_k$ with $PE_{p+k}=R_k\,PE_p$ for all $p$
	(asserted), so relative offsets are linearly accessible; no parameters, and defined for any length.
	\ii \vocab{Learned absolute} embeddings (BERT, GPT-2): a trainable $T_\text{max}\times\dmodel$ table.
	Flexible, but undefined beyond $T_\text{max}$.
	\ii \vocab{RoPE} (rotary): instead of adding anything, rotate each consecutive pair of query and key
	components by an angle proportional to the position. Then $\langle R_m\bq,R_n\bk\rangle$ depends only on
	$m-n$: relative position enters the score directly. Used by many recent LLMs.
	\ii \vocab{ALiBi} \cite{alibi}: no position embeddings at all; add a penalty $-m_h\abs{i-j}$ to the
	attention scores, with a fixed slope $m_h$ per head; for 8 heads the slopes are
	\[ \gv{c08_attention/alibi_slopes}. \]
	Designed for extrapolation to longer sequences than seen in training.
\end{itemize}

\begin{center}
\begin{tikzpicture}
\begin{axis}[width=0.62\linewidth, height=0.42\linewidth, enlargelimits=false, axis on top,
	xlabel={dimension}, ylabel={position}, y dir=reverse, colormap/viridis, colorbar,
	colorbar style={width=0.25cm, font=\scriptsize}, font=\small]
	\addplot[matrix plot*, point meta=explicit, mesh/cols=32] table[x=dim, y=pos, meta=v] {gen/dl/data/c08_attention-pe.dat};
\end{axis}
\end{tikzpicture}
\end{center}

\codefile[\script{c08_attention}]{gen/dl/snip/c08_attention-rope.py}

The script checks that $\langle\mathrm{rope}(\bq,m),\mathrm{rope}(\bk,n)\rangle$ is the same
($\gv{c08_attention/rope_dot}$) for $(m,n)=(3,1)$, $(10,8)$, $(52,50)$, and that the rotation preserves norms.

\section{The Transformer block}
\prototype{One encoder layer: $\X\gets\X+\MHA(\LN(\X))$, then $\X\gets\X+\FFN(\LN(\X))$.}

\subsection{Sublayers}
\begin{itemize}
	\ii \textbf{Self-attention} mixes information \emph{across positions}.
	\ii \textbf{Position-wise feed-forward network}: the same two-layer MLP applied to every position
	independently, $\FFN(\x)=\W_2\,\phi(\W_1\x+\bb_1)+\bb_2$ with hidden size $\dff$ (typically
	$4\dmodel$) and $\phi$ = ReLU or GELU. It mixes information \emph{across features}; most of the
	parameters live here.
	\ii \textbf{Residual connections} around each sublayer (gradient highway, as in ResNet) and
	\textbf{LayerNorm} (\cref{ch:reg}).
\end{itemize}
The original Transformer is \vocab{post-LN}: $\X\gets\LN(\X+\mathrm{Sublayer}(\X))$. Most modern models
are \vocab{pre-LN}: $\X\gets\X+\mathrm{Sublayer}(\LN(\X))$, with a final LayerNorm at the end of the stack.
Pre-LN keeps an un-normalised identity path from input to output, trains more stably at depth, and is
far less sensitive to warm-up; post-LN can reach slightly better results when it trains.

\subsection{Encoder and decoder}
\begin{center}
\begin{tikzpicture}[node distance=0.32cm, every node/.style={font=\small}]
	\node[blockbox, fill=gray!8] (eemb) {token embedding + positions};
	\node[blockbox, fill=orange!15, above=of eemb] (esa) {multi-head self-attention};
	\node[blockbox, fill=yellow!15, above=of esa] (ean) {add \& LayerNorm};
	\node[blockbox, fill=blue!10, above=of ean] (effn) {feed-forward};
	\node[blockbox, fill=yellow!15, above=of effn] (ean2) {add \& LayerNorm};
	\node[draw, dashed, rounded corners, fit=(esa)(ean2), inner sep=4pt, label=left:{$N\times$}] (encbox) {};
	\node[above=0.4cm of ean2] (encout) {encoder output (memory)};
	\draw[flow] (eemb) -- (esa); \draw[flow] (esa) -- (ean); \draw[flow] (ean) -- (effn);
	\draw[flow] (effn) -- (ean2); \draw[flow] (ean2) -- (encout);
	\node[blockbox, fill=gray!8, right=2.6cm of eemb] (demb) {shifted target embedding + positions};
	\node[blockbox, fill=orange!15, above=of demb] (dsa) {masked self-attention};
	\node[blockbox, fill=yellow!15, above=of dsa] (dan) {add \& LayerNorm};
	\node[blockbox, fill=red!10, above=of dan] (dca) {cross-attention (K, V from encoder)};
	\node[blockbox, fill=yellow!15, above=of dca] (dan2) {add \& LayerNorm};
	\node[blockbox, fill=blue!10, above=of dan2] (dffn) {feed-forward};
	\node[blockbox, fill=yellow!15, above=of dffn] (dan3) {add \& LayerNorm};
	\node[draw, dashed, rounded corners, fit=(dsa)(dan3), inner sep=4pt, label=right:{$N\times$}] {};
	\node[blockbox, fill=gray!8, above=0.4cm of dan3] (lin) {linear (often tied) + softmax};
	\draw[flow] (demb) -- (dsa); \draw[flow] (dsa) -- (dan); \draw[flow] (dan) -- (dca);
	\draw[flow] (dca) -- (dan2); \draw[flow] (dan2) -- (dffn); \draw[flow] (dffn) -- (dan3); \draw[flow] (dan3) -- (lin);
	\draw[flow] (encout.east) -| ($(dca.west)+(-0.5,0)$) -- (dca.west);
\end{tikzpicture}
\end{center}
(Residual arrows around each sublayer are omitted; the diagram shows the post-LN ordering of the
original paper.) The decoder has three sublayers: \emph{masked} self-attention over the target prefix,
\vocab{cross-attention} whose queries come from the decoder and whose keys and values come from the
encoder output, and the FFN. Cross-attention shapes: queries $\shp{B,T_t,\dmodel}$, keys and values
$\shp{B,T_s,\dmodel}$, scores $\shp{B,h,T_t,T_s}$, output $\shp{B,T_t,\dmodel}$. The output layer
often shares its weight matrix with the input embedding (\vocab{weight tying}).

Three families use these blocks: \vocab{encoder-only} (BERT; bidirectional attention, for
understanding tasks), \vocab{decoder-only} (GPT-style; causal self-attention only, for generation) and
\vocab{encoder--decoder} (original Transformer, T5; for sequence-to-sequence). \Cref{ch:llm} picks this up.

\section{Counting parameters and FLOPs}
\prototype{One encoder layer with $\dmodel=512$, $\dff=2048$: $\gv{c08_attention/attn512}$ in attention
plus $\gv{c08_attention/ffn512}$ in the FFN, $\gv{c08_attention/enc512}$ with the two LayerNorms.}

\codefile[\script{c08_attention}]{gen/dl/snip/c08_attention-count.py}

\begin{itemize}
	\ii Encoder layer: $(4\dmodel^2+4\dmodel)+(2\dmodel\dff+\dff+\dmodel)+4\dmodel$. With $\dff=4\dmodel$
	this is about $12\dmodel^2$: a third in attention, two thirds in the FFN.
	\ii Decoder layer: an extra attention block and LayerNorm: $\gv{c08_attention/dec512}$ for $\dmodel=512$.
	\ii Embeddings: $V\dmodel$ (plus $T_\text{max}\dmodel$ for learned positions).
\end{itemize}
For $\dmodel=512$ and $768$ the script asserts both layer formulas against PyTorch's
encoder and decoder layer classes.

\begin{example}[Whole models, counted]
	Three configurations, counted by the script:
	\begin{itemize}
		\ii \pt{nn.Transformer(512, 8, 6, 6, 2048)} (6 encoder and 6 decoder layers, the configuration of the
		original ``base'' model): $\gv{c08_attention/tf_body_exact}$ parameters in the stack, asserted. With a
		shared $37{,}000$-token embedding (an illustrative vocabulary size), $\gv{c08_attention/tf_emb}$ more,
		about $\gv{c08_attention/tf_total_shared}$ in total.
		\ii A BERT-base-shaped encoder ($L=12$, $\dmodel=768$, $\dff=3072$, $V=30{,}522$, 512 learned positions,
		2 segment embeddings, a pooler): $\gv{c08_attention/bert_layer}$ per layer, $\gv{c08_attention/bert_body}$
		in the 12 layers, $\gv{c08_attention/bert_emb}$ in embeddings, $\gv{c08_attention/bert_total}$ in total.
		\ii A GPT-2-small-shaped decoder-only model ($L=12$, $\dmodel=768$, $V=50{,}257$, 1024 positions, tied
		output): $\gv{c08_attention/gpt_total}$, of which $\gv{c08_attention/gpt_emb_frac}\%$ is embeddings.
	\end{itemize}
	These are counts of the configurations as defined here, computed by the script, not quotations of
	anyone's published figure; small differences from published numbers come from details such as biases
	and output heads. Rule of thumb for large models: $\approx12L\dmodel^2$ ($\gv{c08_attention/twelve_d2}$ per
	layer at $\dmodel=768$) plus embeddings.
\end{example}

\subsection{Complexity}
Per layer and sequence of length $n$: the projections and FFN cost $O(n\dmodel^2)$; the score matrix
$\Q\K\T$ and the product with $\V$ cost $O(n^2\dmodel)$ and need $O(n^2)$ memory per head. For
$\dmodel=768$ (FLOPs, one layer, one sequence):
\begin{center}
	\small
	\begin{tabular}{rccc}
		\toprule
		$n$ & attention $4n^2\dmodel$ & projections $8n\dmodel^2$ & attention probabilities, 12 heads, fp16 \\
		\midrule
		512 & $\gv{c08_attention/attn_flops_512}$ & $\gv{c08_attention/proj_flops_512}$ & $\gv{c08_attention/attnmat_512}$B \\
		4{,}096 & $\gv{c08_attention/attn_flops_4096}$ & $\gv{c08_attention/proj_flops_4096}$ & $\gv{c08_attention/attnmat_4096}$B \\
		32{,}768 & $\gv{c08_attention/attn_flops_32768}$ & $\gv{c08_attention/proj_flops_32768}$ & $\gv{c08_attention/attnmat_32768}$B \\
		\bottomrule
	\end{tabular}
\end{center}
Attention overtakes the projections once $n>2\dmodel$. The quadratic term is the motivation for sparse,
linear and windowed attention variants, and for FlashAttention's memory savings.

\subsection{KV cache, MQA and GQA}
A decoder generates one token at a time. Without care, step $t$ recomputes keys and values for all $t$
previous tokens in every layer. Since those never change (causal masking), they are stored: the
\vocab{KV cache} holds $\K$ and $\V$ for every layer, head and past position, and each new token computes
only its own $\bq,\bk,\bv$ and attends to the cache. Memory:
\[ 2\times L\times n_\text{kv heads}\times d_\text{head}\times T\times\text{bytes per value}\ \text{per sequence}. \]
For an illustrative configuration with $32$ layers, $32$ heads of size $128$, fp16 and $4096$ tokens, that is
$\gv{c08_attention/kv_per_tok}$\,MiB per token and $\gv{c08_attention/kv_bytes}$\,GiB per sequence.
\vocab{Multi-query attention} (one shared K/V head) cuts it to $\gv{c08_attention/kv_mqa}$\,GiB;
\vocab{grouped-query attention} with $8$ K/V heads shared by groups of query heads to
$\gv{c08_attention/kv_gqa8}$\,GiB. The cache is why long-context serving is memory-bound.

\begin{remark}
	\trap ``The KV cache stores the attention weights'' or ``stores the queries''. It stores keys and values of
	past tokens; queries are needed only for the current token, and attention weights are recomputed each step
	because a new query changes them. The cache helps decoding (generation), not training, where all positions
	are processed in parallel anyway.
\end{remark}

\subsection{FlashAttention, in concept}
FlashAttention \cite{flashattn} computes \emph{exact} attention without ever materialising the $n\times n$
matrix in slow GPU memory: it streams blocks of keys and values through fast on-chip memory and maintains,
per query, a running maximum, a running softmax denominator and a running output, rescaling them when a new
block raises the maximum (an \vocab{online softmax}); the backward pass recomputes blocks instead of storing
them. FLOPs are unchanged; memory traffic, and so wall-clock time and memory, drop sharply.

\codefile[\script{c08_attention}]{gen/dl/snip/c08_attention-online.py}

The blockwise version (keys processed four at a time) equals standard attention to machine precision.

\section{Transformer versus RNN and CNN}
\prototype{Relating the first and the five-hundredth token: $499$ recurrent steps in an RNN, about
$\log_k500$ dilated convolution layers, one attention step.}

\begin{center}
	\small
	\begin{tabular}{lccc}
		\toprule
		layer type & cost per layer & sequential steps & longest path between positions \\
		\midrule
		self-attention & $O(n^2d)$ & $O(1)$ & $O(1)$ \\
		recurrent & $O(nd^2)$ & $O(n)$ & $O(n)$ \\
		convolution, kernel $k$ & $O(knd^2)$ & $O(1)$ & $O(\log_kn)$ (dilated) or $O(n/k)$ \\
		\bottomrule
	\end{tabular}
\end{center}
For seq2seq, the Transformer wins on (1) \textbf{parallel training}: no sequential dependence across time
during training, so GPUs are saturated; (2) \textbf{short paths}: any two positions interact directly, so
long-range dependencies are learned easily and gradients do not pass through hundreds of steps; (3)
\textbf{content-based, dynamic weights}: attention weights depend on the input, unlike a convolution's fixed
kernel; (4) \textbf{scaling}: performance keeps improving with data and parameters. It loses on cost at very
long $n$ and needs positional encodings and lots of data; RNNs remain attractive for streaming with constant
memory per step.

\section{Capstone: an encoder--decoder Transformer, twice}
\prototype{Train on ``reverse this digit sequence'' (lengths 3--8, padded), reach
$\gv{c08_capstone/acc_test}\%$ exact-match accuracy, and reproduce every logit in NumPy.}

\subsection{The PyTorch model}
Pre-LN encoder and decoder layers built from scratch (no \pt{nn.Transformer}), $\dmodel=64$, $h=4$,
$\dff=128$, two encoder and two decoder layers, sinusoidal positions, embeddings scaled by $\sqrt{\dmodel}$ and
shared between source, target and the output layer: $\gv{c08_capstone/n_params}$ parameters. Vocabulary: PAD,
BOS, EOS and ten digits.

\codefile[\script{c08_capstone}]{gen/dl/snip/c08_capstone-torch_model.py}

Before training, the script copies weights into PyTorch's own modules and asserts agreement:
our \pt{MHA} against \pt{nn.MultiheadAttention} (cross-attention with key padding), our encoder layer
against \pt{nn.TransformerEncoderLayer(norm\_first=True)} and our decoder layer against
\pt{nn.TransformerDecoderLayer(norm\_first=True)} (causal and padding masks).

\subsection{Training}
\codefile[\script{c08_capstone}]{gen/dl/snip/c08_capstone-train.py}
AdamW with linear warm-up and cosine decay, gradient clipping at norm 1, batches of 64 random sequences,
$\gv{c08_capstone/steps}$ steps (seconds on a CPU). The decoder input is the target shifted right (starting
with BOS) and the loss ignores PAD positions. The loss falls from $\gv{c08_capstone/loss_first}$ to
$\gv{c08_capstone/loss_last}$ (mean of the last 50 steps).
\begin{center}
\begin{tikzpicture}
\begin{axis}[napkinplot, width=0.6\linewidth, height=0.34\linewidth, xlabel={step}, ylabel={loss}, ymode=log]
	\addplot table[x=step,y=loss] {gen/dl/data/c08_capstone-loss.dat};
\end{axis}
\end{tikzpicture}
\end{center}

\subsection{Decoding and results}
\codefile[\script{c08_capstone}]{gen/dl/snip/c08_capstone-greedy.py}
On $1000$ fresh sequences greedy decoding gives $\gv{c08_capstone/acc_test}\%$ exact-match accuracy
($\gv{c08_capstone/tok_acc}\%$ of tokens). Example: input \pt{\gv{c08_capstone/ex_src}}, output
\pt{\gv{c08_capstone/ex_pred}}. The cross-attention of the last decoder layer (averaged over heads) for the
input \pt{\gv{c08_capstone/x_src}} shows what the model learned: output position $t$ attends to input position
$7-t$ (counting from $0$), the anti-diagonal (average weight there $\gv{c08_capstone/anti_diag}$); the last row is the EOS step.
\begin{center}
\begin{tikzpicture}
\begin{axis}[width=0.5\linewidth, height=0.52\linewidth, enlargelimits=false, axis on top,
	xlabel={source position}, ylabel={output position}, y dir=reverse, colormap/viridis, colorbar,
	colorbar style={width=0.25cm, font=\scriptsize}, point meta min=0, point meta max=1, font=\small,
	xtick={0,...,7}, ytick={0,...,8}]
	\addplot[matrix plot*, point meta=explicit, mesh/cols=8] table[x=src, y=tgt, meta=a] {gen/dl/data/c08_capstone-xattn.dat};
\end{axis}
\end{tikzpicture}
\end{center}

\subsection{The NumPy port}
The same forward pass in NumPy, reading the trained \pt{state\_dict}:
\codefile[\script{c08_capstone}]{gen/dl/snip/c08_capstone-numpy_model.py}
In float64 the NumPy logits equal PyTorch's on 50 test sequences (to $10^{-9}$), and greedy decoding in NumPy
produces exactly the same 200 output sequences as PyTorch ($\gv{c08_capstone/np_acc200}\%$ correct). Writing the
NumPy version is the best interview preparation in this book: every reshape, transpose and mask has to be right
for the assertion to pass.

\begin{moral}
	Attention is a softmax-weighted average with learned queries, keys and values, scaled by $\sqrt{d_k}$ to keep
	the softmax soft. A Transformer layer is attention (mix positions) plus an MLP (mix features), each wrapped
	in a residual connection and LayerNorm; positions must be injected; cost is quadratic in length.
\end{moral}

\section\problemhead

\begin{problem}\onechili
	Two queries $\bq_1=(1,0)$, $\bq_2=(0,2)$, keys $\bk_1=(1,1)$, $\bk_2=(0,1)$, values $v_1=10$, $v_2=20$,
	$d_k=2$. Compute both attention outputs.
	\begin{hint}Scores $\bq\T\bk/\sqrt2$, softmax per query.\end{hint}
	\begin{sol}
		Query 1: scores $(1,0)/\sqrt2=(0.707,0)$, weights $(\gv{c08_attention/pa_w1})$, output
		$\gv{c08_attention/pa_o1}$. Query 2: scores $(2,2)/\sqrt2$, weights $(\gv{c08_attention/pa_w2})$,
		output $\gv{c08_attention/pa_o2}$ (checked against PyTorch).
	\end{sol}
\end{problem}

\begin{problem}\onechili
	Give the parameter count of one multi-head self-attention block with $\dmodel=256$ and $h=4$, and of
	the FFN with $\dff=1024$, biases included.
	\begin{hint}$4d^2+4d$ and $2d\dff+\dff+d$.\end{hint}
	\begin{sol}Attention $4\cdot65{,}536+1{,}024=\gv{c08_attention/pb_attn}$; FFN
		$524{,}288+1{,}024+256=\gv{c08_attention/pb_ffn}$.\end{sol}
\end{problem}

\begin{problem}\twochili
	Scores for one query over four keys are $(8,4,0,-4)$ with $d_k=64$ but someone forgot the scaling.
	Compare the softmax weights with and without dividing by $\sqrt{d_k}$.
	\begin{hint}$\sqrt{64}=8$.\end{hint}
	\begin{sol}
		Unscaled: $\softmax(8,4,0,-4)=(\gv{c08_attention/pc_raw})$, essentially one-hot, so almost no gradient
		reaches keys 2--4. Scaled: $\softmax(1,0.5,0,-0.5)=(\gv{c08_attention/pc_scaled})$, a soft distribution
		that still prefers key 1.
	\end{sol}
\end{problem}

\begin{problem}\twochili
	What are the shapes of $\Q$, $\K$, $\V$, the score tensor and the output in the cross-attention of a decoder
	with batch $B=2$, target length $5$, source length $9$, $\dmodel=32$, $h=4$?
	\begin{hint}Queries come from the decoder.\end{hint}
	\begin{sol}
		After splitting heads: $\Q$ $\shp{2,4,5,8}$, $\K$ and $\V$ $\shp{2,4,9,8}$, scores $\shp{2,4,5,9}$, per-head
		output $\shp{2,4,5,8}$, merged output $\shp{2,5,32}$.
	\end{sol}
\end{problem}

\begin{problem}\onechili
	Why does a decoder-only language model need a causal mask during training but an encoder like BERT does not?
	\begin{hint}What is each trained to predict?\end{hint}
	\begin{sol}
		The decoder predicts token $t+1$ from tokens $\le t$ at every position simultaneously; without the mask
		position $t$ could read token $t+1$ and the loss would be trivial. BERT predicts masked-out tokens from
		both sides; the information to hide is removed from the input (the mask token), so all positions may attend
		to all others.
	\end{sol}
\end{problem}

\begin{problem}\twochili
	Estimate the KV-cache size for a model with $24$ layers, $16$ heads of size $64$, fp16, a batch of $8$
	sequences of $2048$ tokens. How does grouped-query attention with $4$ K/V heads change it?
	\begin{hint}$2\cdot L\cdot n_\text{kv}\cdot d_\text{head}\cdot T\cdot B\cdot2$ bytes.\end{hint}
	\begin{sol}
		$2\cdot24\cdot16\cdot64\cdot2048\cdot8\cdot2$ bytes $=\gv{c08_attention/pf_kv}$\,GiB. With $4$ K/V heads it
		is a quarter, $\gv{c08_attention/pf_kv4}$\,GiB.
	\end{sol}
\end{problem}
